\section{问题一模型的建立与求解}

\subsection{问题一的描述与分析}

问题一需要完成“统计分析—短期预测—基础调度”三个连续任务：先从附件中识别不同区域、不同任务类型的 GPU 需求结构，再建立未来 24\,h 的短期需求预测模型，最后对 2376--2399\,h 实际到达任务进行基础算力调度，并在 2400--2405\,h 收尾时域内结清跨越 2399\,h 的任务，给出任务调度甘特图和区域 GPU 利用率。

本问的预测时域仅为 24\,h，因此不追求复杂长期时序结构，而先检验“小时到达是否近似 Poisson、区域与任务类型是否独立、资源需求应按何种条件分层”。数据诊断结果分别决定计数过程、业务标记和资源标记，形成结构化标记复合 Poisson 模型；随后在不引入电价、碳强度和储能决策的条件下建立本地优先、最小调整的基础调度模型，为后续两问提供可解释的计算侧基准。总体流程见图\ref{fig:q1-flow}。

\begin{figure}[!htbp]
\centering
\includegraphics[width=\FigureFlowWidth]{\ProjectRoot/modules/20_q1/figures/q1_flow.png}
\caption{问题一建模求解总体流程}
\label{fig:q1-flow}
\end{figure}

\subsection{数据预处理与统计分析}

为避免测试信息反向影响模型结构，\textbf{所有用于确定预测模型形式的统计诊断均以 0--2351\,h 训练段为准}；若使用全 50000 个任务的统计量，仅用于描述整体业务规模和结果解释，不参与模型选择、参数回调或测试集再拟合。

\subsubsection{小时到达过程诊断}

基于 0--2351\,h 训练段，小时任务到达数均值为 $\hat\lambda\approx20.8$，Fano 因子约为 0.97，接近 Poisson 计数过程的方差均值比 1。经验频率与 Poisson 理论概率整体吻合；自相关函数中 lag~1、24、168 均接近 0，未形成稳定的短期记忆、日周期或周周期结构。图\ref{fig:q1-arrival-diagnosis}给出相应诊断，其中 $\pm0.0404$ 为白噪声近似下的 95\% 参考带。

因此，在本题 24\,h 短期预测尺度内采用齐次 Poisson 到达过程即可保留数据支持的主要结构，没有必要额外引入 AR、季节项或黑箱时序模型。该结论仅针对本题短预测窗，不解释为长期到达机制恒定。

\begin{figure}[!htbp]
\centering
\includegraphics[width=\FigureDoubleWidth]{\ProjectRoot/modules/20_q1/figures/fig_q1_poisson_diagnostic.png}
\caption{小时到达任务数的 Poisson 拟合与自相关诊断}
\label{fig:q1-arrival-diagnosis}
\end{figure}
\FloatBarrier

\subsubsection{区域与任务类型的联合结构}

任务类型（TaskType）与来源区域（SourceRegion）并非独立。训练段统计得到 Cram\'{e}r's~V$\approx0.434$：实时推理任务主要集中于 RegionA/B/C，AI 训练任务在 RegionD/E/F 中更为集中，而批量推理相对均匀。因此
\begin{equation}
P(R,K)\neq P(R)P(K),
\label{eq:q1-rk-nonindependent}
\end{equation}
业务组成必须用区域—类型联合概率描述，不能分别估计两个边际分布后相乘。图\ref{fig:q1-composition}展示了该结构。

从计算工作量看，AI 训练任务数约占全部任务的三分之一，却贡献超过 80\% 的 GPU-hour。任务数与计算负荷强度显著错位，说明预测和调度均需保留 GPU 资源标记，不能只预测任务到达数。

\begin{figure}[!htbp]
\centering
\includegraphics[width=\FigureDoubleWidth]{\ProjectRoot/modules/20_q1/figures/fig_q1_task_structure.png}
\caption{区域—任务类型结构及其计算工作量差异}
\label{fig:q1-composition}
\end{figure}
\FloatBarrier

\subsubsection{条件资源结构}

三类任务的 GPU 需求具有稳定分层：实时推理中位数约 4 GPU，批量推理居中，AI 训练中位数约 72 GPU。给定 TaskType 后，SourceRegion 对 GPU 需求和执行时长的附加解释力较弱；GPU 与执行时长的线性相关接近 0，但低线性相关并不能推出统计独立。

据此，本文保留
\begin{equation}
(G,D)\mid K=k\sim\hat F_{GD}^{(k)}
\label{eq:q1-resource-structure}
\end{equation}
的联合经验分布。短期 GPU 总量预测实际使用该联合分布诱导的 GPU 边际；执行时长 $D$ 则用于保持任务资源结构与基础调度的一致性。

综合以上三项诊断，问题一的预测模型结构由数据直接确定为
\[
\boxed{\text{Poisson 计数过程}+\text{区域—类型联合业务标记}+\text{按任务类型条件化的资源标记}}.
\]

\subsection{结构化标记复合 Poisson 预测模型}

\subsubsection{到达过程与联合业务标记}

设第 $t$ 小时到达任务总数为 $N_t$，采用
\begin{equation}
N_t\sim\mathrm{Poisson}(\lambda).
\label{eq:q1-arrival}
\end{equation}
每个任务携带来源区域 $R$ 与任务类型 $K$，联合概率为
\begin{equation}
P(R=r,K=k)=\pi_{rk},\qquad \sum_{r=1}^{6}\sum_{k=1}^{3}\pi_{rk}=1.
\label{eq:q1-pi}
\end{equation}
由 Poisson 分裂性质，区域—任务类型子流满足
\begin{equation}
N_{rkt}\sim\mathrm{Poisson}(\lambda\pi_{rk}).
\label{eq:q1-split}
\end{equation}

\subsubsection{资源标记与 GPU 需求}

给定任务类型 $K=k$，GPU 需求 $G$ 与执行时长 $D$ 从训练样本的联合经验分布中抽取：
\begin{equation}
(G,D)\mid K=k\sim\hat F_{GD}^{(k)}.
\label{eq:q1-mark}
\end{equation}
因此完整任务标记关系可写为
\begin{equation}
P(R,K,G,D)=\pi_{rk}\hat f_{GD}^{(k)}(g,d).
\label{eq:q1-joint}
\end{equation}

定义第 $t$ 小时 Region $r$、TaskType $k$ 的新到达 GPU 需求为
\begin{equation}
Y_{rkt}=\sum_{i:R_i=r,K_i=k,A_i=t}G_i.
\label{eq:q1-demand}
\end{equation}
则复合 Poisson 性质给出
\begin{equation}
\mathbb E[Y_{rkt}]=\lambda\pi_{rk}\mathbb E[G\mid K=k],\qquad
\mathrm{Var}(Y_{rkt})=\lambda\pi_{rk}\mathbb E[G^2\mid K=k].
\label{eq:q1-demand-moments}
\end{equation}
系统层、区域层和任务类型层预测均由 $Y_{rkt}$ 加总得到，从而保持层级严格一致。

\subsubsection{模型汇总}

问题一短期预测模型汇总为
\begin{equation}
\left\{
\begin{aligned}
&N_t\sim\mathrm{Poisson}(\lambda),\\
&P(R=r,K=k)=\pi_{rk},\\
&(G,D)\mid K=k\sim\hat F_{GD}^{(k)},\\
&N_{rkt}\sim\mathrm{Poisson}(\lambda\pi_{rk}),\\
&Y_{rkt}=\sum_{i:R_i=r,K_i=k,A_i=t}G_i,\\
&\mathbb E[Y_{rkt}]=\lambda\pi_{rk}\mathbb E[G\mid K=k].
\end{aligned}
\right.
\label{eq:q1-model-summary}
\end{equation}

\subsection{模型求解与概率预测}

按时间顺序划分训练、验证与测试数据：0--2351\,h 为训练段；2352--2375\,h 为独立验证段，用于检查预测区间覆盖和模型稳定性，\textbf{不依据验证或测试表现切换模型结构}；随后在 0--2375\,h 上重新估计最终参数；2376--2399\,h 作为一次性测试段，不用于回调参数。

训练段参数采用频率估计：
\begin{equation}
\hat\lambda=\frac1T\sum_{t=0}^{T-1}n_t,\qquad
\hat\pi_{rk}=\frac{c_{rk}}{\sum_{r,k}c_{rk}},\qquad T=2352.
\label{eq:q1-estimation}
\end{equation}
条件资源分布 $\hat F_{GD}^{(k)}$ 直接由对应任务类型的训练样本经验频率得到。

均值预测为
\begin{equation}
\hat Y_{rkt}=\hat\lambda\hat\pi_{rk}\frac1{|\mathcal D_k|}\sum_{i\in\mathcal D_k}G_i.
\label{eq:q1-pred-mean}
\end{equation}
90\% 预测区间采用 Monte Carlo 生成：先抽取 Poisson 小时任务数，再依次抽取联合业务标记和条件资源标记，取模拟分布的 5\% 与 95\% 分位数。具体流程见算法\ref{alg:q1-prediction}。

已有工作负载研究说明，Poisson 近似应由具体时间尺度上的离散度与相关性诊断决定\cite{zou_poissonity_2022}；分层需求预测则强调在细粒度结构与参数稳定性之间进行数据驱动聚合\cite{cohen_dac_2022}。本文据此保留训练段识别出的联合业务标记和类型条件资源标记。

\begin{algorithm}[H]
\SetAlgoLined
\setstretch{0.82}
\SetAlFnt{\footnotesize}
\SetArgSty{textnormal}
\SetKwInput{KwIn}{输入}
\SetKwInput{KwOut}{输出}
\SetKwFor{ForEach}{对每个}{执行}{结束}
\SetKwIF{If}{ElseIf}{Else}{若}{则}{否则若}{否则}{结束}
\SetKwFunction{Diag}{结构诊断}
\SetKwFunction{Fit}{拟合MCP}
\SetKwFunction{MC}{蒙特卡洛预测}
\SetKwFunction{Place}{本地优先放置}
\caption{两阶段短期概率预测与基础调度}
\label{alg:q1-unified}
\KwIn{$\mathcal D_{0:2399}$，$\mathcal J_{2376:2399}$，容量与 SLA，$M=5000$}
\KwOut{24 h 概率预测、实际任务排程与约束审计}

\tcp*[l]{阶段 I：结构识别与独立验证}
$(\mathcal D_{tr},\mathcal D_{val},\mathcal D_{te})\leftarrow(0{:}2351,2352{:}2375,2376{:}2399)$\;
$\mathcal M\leftarrow\Diag(\mathcal D_{tr})$ \tcp*[r]{Fano/ACF、区域--类型关联、$(G,D)\mid K$}
$\theta_1=(\hat\lambda,\hat\pi_{rk},\widehat F_{GD}^{(k)})\leftarrow\Fit(\mathcal M,\mathcal D_{tr})$\;
$\mathcal I_{val}\leftarrow\MC(\theta_1,M)$，计算 MAE、WAPE、RMSE 与 PICP\;
验证通过后冻结 $\mathcal M$，测试段不参与模型选择\;

\tcp*[l]{阶段 II：参数重估、最终预测与基础调度}
$\theta_2\leftarrow\Fit(\mathcal M,\mathcal D_{0:2375})$，$\mathcal I_{te}\leftarrow\MC(\theta_2,M)$\;
初始化 2376 h 前未完成任务形成的 GPU/IT/Facility 占用\;
\ForEach{按到达顺序的 $j\in\mathcal J_{2376:2399}$}{
  构造合法域 $C_j=\{(r,s):\mathrm{SLA},\mathrm{LatestFinish},s+p_j\le2406\}$；RT 固定 $s=a_j$\;
  $(r_j^*,s_j^*)\leftarrow\Place(C_j)$，同时满足 GPU/IT/Facility 容量\;
  \If{$(r_j^*,s_j^*)=\varnothing$}{返回不可行排程\;}
  更新区间 $[s_j^*,s_j^*+p_j)$ 的资源占用\;
}
审计唯一执行、RT 即时开工、容量、SLA、LatestFinish 与 2406 边界\;
\end{algorithm}


\subsection{基础算力调度模型}

\subsubsection{本地优先最小调整原则}

基础调度以“来源区域、到达即执行”为参考状态，仅在硬约束冲突时做必要调整。该参考状态能够在浮点误差范围内复现附件给定的基准 AI IT 负荷，最大误差约 $3.33\times10^{-7}$ MW。问题一刻意不引入电价、碳强度、新能源和储能信号，使本问只回答基础算力可行性问题。

具体规则为：实时推理任务必须到达即开工，源区域不可行时才在满足 20 ms SLA 的区域中迁移；批量推理和 AI 训练允许等待，优先在源区域完整合法时间域中寻找最早可行时刻，只有源区域全域不可行时才迁移到其他 SLA 合法区域。所有任务不可拆分、不可抢占，执行期间区域保持不变，并须满足 LatestFinish 与 2406\,h 终端边界。

\subsubsection{资源占用与 GPU 利用率}

分钟级持续时间换算为小时后，任务 $i$ 在小时 $t$ 的实际重叠量为
\begin{equation}
\omega(i,t)=\max\bigl(0,\min(t+1,s_i+p_i)-\max(t,s_i)\bigr),
\label{eq:q1-overlap}
\end{equation}
并满足 $\sum_t\omega(i,t)=p_i$。区域 GPU 利用率定义为
\begin{equation}
U_{rt}=\frac{CarryInGPU_{rt}+\sum_iG_i\omega(i,t)}{AvailableGPU_r},
\label{eq:q1-utilization}
\end{equation}
其中 $CarryInGPU_{rt}$ 表示 2376\,h 前已开始、但与 $[2376,2406)$ 仍有重叠的历史任务占用。完整调度流程见算法\ref{alg:q1-scheduling}。

\subsection{结果与分析}

\subsubsection{任务结构统计}

表\ref{tab:q1-stat-summary}列出部分代表性区域—任务类型统计。AI 训练在 D/E/F 区域更集中，且平均 GPU 需求约 71 GPU；其 GPU-hour 占比远高于任务数占比。该结果与前述联合业务标记和条件资源标记一一对应，说明统计分析直接参与了模型结构选择。

\begin{table}[!htbp]
\centering
\caption{部分区域—任务类型统计摘要}
\label{tab:q1-stat-summary}
\begin{tabular}{lcccc}
\toprule
\rowcolor{CumcmLightBlue}
区域—任务类型 & 任务数 & GPU-hour & 平均 GPU & 平均时长/min \\
\midrule
RegionA, AITraining & 879 & 214319 & 71.75 & 206.2 \\
RegionA, BatchInference & 3331 & 154341 & 13.53 & 205.6 \\
RegionA, RealTimeInference & 5852 & 79792 & 4.03 & 204.0 \\
RegionD, AITraining & 5726 & 1387123 & 71.20 & 204.5 \\
RegionD, BatchInference & 3346 & 155177 & 13.54 & 204.8 \\
RegionD, RealTimeInference & 488 & 6559 & 4.08 & 192.7 \\
RegionF, AITraining & 4177 & 1032935 & 70.90 & 207.5 \\
RegionF, BatchInference & 2033 & 90263 & 13.37 & 201.1 \\
RegionF, RealTimeInference & 497 & 6837 & 4.02 & 205.9 \\
\bottomrule
\end{tabular}
\end{table}
\FloatBarrier

\subsubsection{24 h 短期预测结果}

测试段预测结果见表\ref{tab:q1-forecast-metrics}。系统层 WAPE 为 0.243，与历史均值基线接近；本模型的主要增量价值在于同时给出区域—任务类型层级一致的概率预测和区间，而不是通过增加短期时序复杂度获得有限的点预测改善。

\begin{table}[!htbp]
\centering
\caption{2376--2399 h 独立测试集预测结果}
\label{tab:q1-forecast-metrics}
\begin{tabular}{lccccc}
\toprule
\rowcolor{CumcmLightBlue}
层级 & MAE & RMSE & WAPE & 90\% PICP & MPIW \\
\midrule
System & 176.68 & 221.05 & 0.243 & 87.50\% & 695.05 \\
Region & 65.49 & 95.00 & 0.540 & 89.58\% & 251.85 \\
TaskType & 67.09 & 122.55 & 0.277 & 90.28\% & 281.03 \\
Region$\times$TaskType & 26.49 & 53.74 & 0.656 & 94.21\% & 97.67 \\
\bottomrule
\end{tabular}
\end{table}
\FloatBarrier

系统层 90\% 区间覆盖率为 87.5\%（21/24），Region 和 TaskType 层分别为 89.58\% 和 90.28\%，Region$\times$TaskType 层为 94.21\%。测试集只有 24 个小时，系统层与名义覆盖率的小幅偏差不作为再次调参依据。图\ref{fig:q1-fig08}给出系统层预测均值、90\% 区间与实际观测值。

\begin{figure}[!htbp]
\centering
\includegraphics[width=\FigureDoubleWidth]{\ProjectRoot/modules/20_q1/figures/fig_q1_forecast_24h.png}
\caption{2376--2399 h 系统层 GPU 需求概率预测}
\label{fig:q1-fig08}
\end{figure}
\FloatBarrier

\subsubsection{基础调度结果}

2376--2399\,h 共 538 个实际到达任务全部完成调度，其中无跨区域迁移任务，仅 4 个弹性任务因局部容量冲突在源区域顺延。最大区域 GPU 利用率为 99.02\%，平均利用率为 37.23\%，最晚任务于 2405.6\,h 完成，未越过 2406\,h 边界。结果见表\ref{tab:q1-schedule-summary}。

\begin{table}[!htbp]
\centering
\caption{基础算力调度结果}
\label{tab:q1-schedule-summary}
\begin{tabular}{lccccc}
\toprule
\rowcolor{CumcmLightBlue}
任务数 & 迁移数 & 顺延数 & 最大 GPU 利用率 & 平均 GPU 利用率 & 最晚完成时间 \\
\midrule
538 & 0 & 4 & 99.02\% & 37.23\% & 2405.6\,h \\
\bottomrule
\end{tabular}
\end{table}
\FloatBarrier

图\ref{fig:q1-schedule}同时给出题目要求的任务甘特图与区域 GPU 利用率。甘特图以 2399\,h 为主调度窗与收尾窗分界：2376--2399\,h 为实际到达窗口，2400--2405\,h 仅用于结清跨越 2399\,h 的在途任务。全部任务保持来源区域执行，4 个顺延任务由本地时间平移消解冲突。

\begin{figure}[!htbp]
\centering
\begin{subfigure}[t]{\FigureDoubleWidth}
\centering
\includegraphics[width=\linewidth]{\ProjectRoot/modules/20_q1/figures/fig_q1_gantt.png}
\caption{2376--2399 h 到达任务及其 2400--2405 h 收尾区间}
\end{subfigure}
\vspace{1em}
\begin{subfigure}[t]{\FigureDoubleWidth}
\centering
\includegraphics[width=\linewidth]{\ProjectRoot/modules/20_q1/figures/fig_q1_gpu_util.png}
\caption{2376--2405 h 六区域 GPU 利用率}
\end{subfigure}
\caption{问题一基础算力调度结果}
\label{fig:q1-schedule}
\end{figure}
\FloatBarrier

\subsection{模型检验}

\subsubsection{到达率敏感性}

在 $\hat\lambda$ 附近施加 $\pm5\%$ 扰动后重新计算系统层预测，WAPE 位于 0.225--0.262，见表\ref{tab:q1-sensitivity}。短期结论对到达率的小幅估计误差保持稳定。

\begin{table}[!htbp]
\centering
\caption{到达率敏感性分析}
\label{tab:q1-sensitivity}
\begin{tabular}{lcccc}
\toprule
\rowcolor{CumcmLightBlue}
情景 & 系统均值 GPU & MAE & RMSE & WAPE \\
\midrule
$\lambda$ 降低 5\% & 582.22 & 190.43 & 238.33 & 0.262 \\
基准经验标记 & 612.86 & 176.68 & 221.05 & 0.243 \\
$\lambda$ 提高 5\% & 643.51 & 163.91 & 206.90 & 0.225 \\
均匀 GPU 标记对照 & 614.29 & 176.09 & 220.32 & 0.242 \\
\bottomrule
\end{tabular}
\end{table}
\FloatBarrier

\subsubsection{区间校准性}

验证段 2352--2375\,h 的系统层 90\% 区间覆盖率为 100\%，MPIW 为 698.05；测试段系统层覆盖率为 87.5\%。Region$\times$TaskType 层验证与测试覆盖率分别为 94.91\% 和 94.21\%。两段结果未显示需要基于测试表现重新调整模型的证据，因此保留训练段诊断所确定的短期预测结构。